% !TeX program = xelatex
% SPDX-License-Identifier: MIT
% Copyright (c) 2026 Ziyu Peng and 刘毅涵
% Author: https://pzy54154631-hub.github.io/homepage/
% Author: https://pzy54154631-hub.github.io/liuyihan/
% 作者：Ziyu Peng、刘毅涵
\documentclass[UTF8,a4paper,fontset=fandol]{ctexart}
\usepackage[margin=2.5cm,headheight=16pt]{geometry}
\usepackage[expansion=false]{microtype}
\usepackage{xcolor,fancyhdr,listings,fvextra,tikz}
\setmainfont{texgyrepagella-regular.otf}[
  BoldFont=texgyrepagella-bold.otf,
  ItalicFont=texgyrepagella-italic.otf,
  BoldItalicFont=texgyrepagella-bolditalic.otf]
\setmonofont{lmmono10-regular.otf}[BoldFont=lmmonolt10-bold.otf,
  ItalicFont=lmmono10-italic.otf]
\definecolor{inkblue}{HTML}{385A75}
\definecolor{paperblue}{HTML}{F1F5F8}
\ctexset{section={format=\Large\bfseries\color{inkblue}}}
\pagestyle{fancy}
\fancyhf{}
\fancyhead[L]{LaTeX 学习手记}
\fancyhead[R]{字体与版面}
\fancyfoot[C]{\thepage}
\lstset{basicstyle=\ttfamily\small,breaklines=true,
  backgroundcolor=\color{paperblue},frame=single,showstringspaces=false}
\usepackage[colorlinks=true,allcolors=inkblue]{hyperref}
\title{让同一种信息拥有同一种样式}
\author{\href{https://pzy54154631-hub.github.io/homepage/}{\textit{Ziyu Peng}}, University of Colorado Boulder
  \\[0.35em]
  \href{https://pzy54154631-hub.github.io/liuyihan/}{刘毅涵}，暨南大学经济学院}
\date{}
\begin{document}
\maketitle
\section{统一视觉层级}
这是一段正文。\textbf{加粗}用于强调关键点，
\textcolor{inkblue}{主题色}帮助识别少量提示信息。
{\kaishu 这一句使用局部楷体。}后面的文字恢复正常。
\section{展示代码}
下面是英文 Python 代码，只作代码排版演示。
\begin{lstlisting}[language=Python]
values = [120, 80, 200]
total = sum(values)
print(total)
\end{lstlisting}
包含中文的 LaTeX 源码用 Verbatim 展示，不把它当作正文执行。
\begin{Verbatim}[breaklines=true,fontsize=\small,frame=single]
\section{材料整理}
这是中文说明。正文中的百分号写作 \%。
\end{Verbatim}
\section{一个小流程图}
\begin{center}
\begin{tikzpicture}[every node/.style={draw=inkblue,rounded corners,
  minimum height=10mm,minimum width=25mm,fill=paperblue}]
  \node (a) at (0,0) {输入材料};
  \node (b) at (4,0) {核对来源};
  \node (c) at (8,0) {输出文档};
  \draw[->,thick,inkblue] (a) -- (b);
  \draw[->,thick,inkblue] (b) -- (c);
\end{tikzpicture}
\end{center}
\end{document}
